% Generated draft: requires mathematical and source review.
\documentclass[11pt]{article}
\usepackage[margin=25mm]{geometry}
\usepackage{fontspec}
\setmainfont{Latin Modern Roman}
\usepackage{amsmath,amssymb,mathtools,unicode-math}
\setmathfont{Latin Modern Math}
\usepackage{enumitem,needspace,xurl,hyperref}
\setlength{\parindent}{0pt}
\setlength{\parskip}{5pt}
\setlength{\emergencystretch}{3em}
\newcommand{\N}{\mathbb{N}}
\newcommand{\Z}{\mathbb{Z}}
\newcommand{\Q}{\mathbb{Q}}
\newcommand{\R}{\mathbb{R}}
\newcommand{\C}{\mathbb{C}}
\newcommand{\F}{\mathbb{F}}
\newcommand{\A}{\mathbb{A}}
\newcommand{\PP}{\mathbb{P}}
\newcommand{\E}{\mathbb{E}}
\newcommand{\eps}{\varepsilon}
\newcommand{\dd}{\,\mathrm{d}}
\begin{document}
\section*{UnsolvedMath ID: 20000305}
\textbf{Source problem code:} catalogue entry
\subsection{Definitions and mathematical statement}
Let $R=\C[[x_1,\ldots,x_n]]$ with maximal ideal $\mathfrak m$, and let $f$ be a nonzero nonunit formal series for the intended finite-threshold question. The original does not explicitly exclude zero or units; their threshold conventions must be stated separately. For a convergent series $g$, the usual local log-canonical threshold at the origin is computed from a log resolution by the minimum of log-discrepancy divided by vanishing order over divisors whose centers contain the origin. The source cites this as rational.

For arbitrary formal $f$, let $f_{\leq N}$ denote its total-degree polynomial truncation. The source defines, using its cited approximation results,
\[
 \operatorname{lct}_0(f)=\lim_{N\to\infty}\operatorname{lct}_0(f_{\leq N})
\]
once the truncations are nonzero, and asks whether this number is rational. The problem concerns every such formal series, not only convergent or finitely determined examples.

The supplied research report attributes an affirmative published answer to de Fernex--Musta\c t\u a: fixed-dimensional threshold sets are closed and formal, polynomial, and geometric hypersurface threshold sets coincide. It also invokes Temkin resolution on regular excellent characteristic-zero formal schemes for a direct divisorial proof. Zero-threshold and set-closure conventions and precise references are not supplied.

Its auxiliary candidate family takes $n\geq2$, positive integral exponents $a_i$, $f_0=\sum_i x_i^{a_i}$, $A=\max_i a_i$, and $\sum_i1/a_i\leq1$. It claims, for every formal $h\in\mathfrak m^{A+1}$,
\[
 \operatorname{lct}_0(f_0+h)=\sum_i\frac1{a_i}.
\]
It calls $A+1$ the optimal uniform total-degree radius, because cancellation by $h=-x_j^A$ for an index with $a_j=A$ lies in $\mathfrak m^A$ and changes the threshold. The report states this in characteristic zero; the exact field/formal-resolution hypotheses require review. Neither the general rationality theorem nor the sharp-family claim is independently proved here.
\subsection{Short English statement}
Is the log-canonical threshold defined by truncation limits of a formal complex power series always rational? The report says this is answered affirmatively in the literature and separately proposes a sharp finite-jet stability family. Missing zero/unit conventions and exact references remain for review.
\subsection{Sources}
\begin{itemize}
\item Local UnsolvedMath registry, permanent ID $20000305$, source code \texttt{catalogue entry}; created and updated 2026-08-14. Status \texttt{solved}; classification \texttt{SOLVED-IN-LITERATURE}; difficulty L3; proposer field null, while verification attributes the original to Mircea Musta\c t\u a.
\item AIM Problem Lists, Algebraic geometry, workshop \emph{Numerical invariants of singularities and higher-dimensional algebraic varieties}, Problem 2, source type PDF, \url{https://aimath.org/WWN/singularvariety/problems_final.pdf}. Canonical location \texttt{aim-algebraic-geometry-notes.json}, \texttt{notes[304]}, tag \texttt{problem}; statement status \texttt{exact}. Verification notes OCR damage in the name, subscripts, and line breaks. Original remarks and literature field empty.
\item Original reference labels [1] and [3] support the analytic definition/rationality, and [1] supports the truncation-limit definition, but their bibliographic entries are absent. The research summary names de Fernex--Musta\c t\u a and Temkin without exact works or identification with those labels.
\item Supplied research attempt 1: result \texttt{partial\_result}, run status \texttt{solved}, rigor label \texttt{proved}, low novelty confidence. Candidate contribution \texttt{sharp\_worked\_family}: stability on $f_0+\mathfrak m^{A+1}$ and sharpness through degree-$A$ cancellation. Explicit caveat: machine-generated research aid, not peer reviewed; independently verify before citation.
\item Dataset attribution: Ulam AI / UnsolvedMath Contributors, CC BY 4.0; original problem attribution remains to the supplied AIM workshop collection.
\end{itemize}
\end{document}
